
\documentclass{article}
\usepackage{colortbl}
\usepackage{makecell}
\usepackage{multirow}
\usepackage{supertabular}

\begin{document}

\newcounter{utterance}

\centering \large Interaction Transcript for game `hot\_air\_balloon', experiment `air\_balloon\_survival\_de\_negotiation\_hard', episode 1 with gpt{-}oss{-}120b{-}t1.0.
\vspace{24pt}

{ \footnotesize  \setcounter{utterance}{1}
\setlength{\tabcolsep}{0pt}
\begin{supertabular}{c@{$\;$}|p{.15\linewidth}@{}p{.15\linewidth}p{.15\linewidth}p{.15\linewidth}p{.15\linewidth}p{.15\linewidth}}
    \# & $\;$A & \multicolumn{4}{c}{Game Master} & $\;\:$B\\
    \hline

    \theutterance \stepcounter{utterance}  
    & & \multicolumn{4}{p{0.6\linewidth}}{
        \cellcolor[rgb]{0.9,0.9,0.9}{
            \makecell[{{p{\linewidth}}}]{
                \texttt{\tiny{[P1$\langle$GM]}}
                \texttt{Du nimmst an einem kollaborativen Verhandlungspiel Teil.} \\
\\ 
\texttt{Zusammen mit einem anderen Teilnehmer musst du dich auf eine Reihe von Gegenständen entscheiden, die behalten werden. Jeder von euch hat eine persönliche Verteilung über die Wichtigkeit der einzelnen Gegenstände. Jeder von euch hat eine eigene Meinung darüber, wie wichtig jeder einzelne Gegenstand ist (Gegenstandswichtigkeit). Du kennst die Wichtigkeitsverteilung des anderen Spielers nicht. Zusätzlich siehst du, wie viel Aufwand jeder Gegenstand verursacht.  } \\
\texttt{Ihr dürft euch nur auf eine Reihe von Gegenständen einigen, wenn der Gesamtaufwand der ausgewählten Gegenstände den Maximalaufwand nicht überschreitet:} \\
\\ 
\texttt{Maximalaufwand = 4026} \\
\\ 
\texttt{Hier sind die einzelnen Aufwände der Gegenstände:} \\
\\ 
\texttt{Aufwand der Gegenstände = \{"C33": 780, "A89": 742, "A48": 997, "A04": 665, "A41": 140, "C64": 77, "C49": 513, "A61": 383, "C04": 587, "B53": 825, "B85": 319, "B61": 448, "A45": 516, "A03": 694, "B21": 366\}} \\
\\ 
\texttt{Hier ist deine persönliche Verteilung der Wichtigkeit der einzelnen Gegenstände:} \\
\\ 
\texttt{Werte der Gegenstandswichtigkeit = \{"C33": 138, "A89": 583, "A48": 868, "A04": 822, "A41": 783, "C64": 65, "C49": 262, "A61": 121, "C04": 508, "B53": 780, "B85": 461, "B61": 484, "A45": 668, "A03": 389, "B21": 808\}} \\
\\ 
\texttt{Ziel:} \\
\\ 
\texttt{Dein Ziel ist es, eine Reihe von Gegenständen auszuhandeln, die dir möglichst viel bringt (d. h. Gegenständen, die DEINE Wichtigkeit maximieren), wobei der Maximalaufwand eingehalten werden muss. Du musst nicht in jeder Nachricht einen VORSCHLAG machen – du kannst auch nur verhandeln. Alle Taktiken sind erlaubt!} \\
\\ 
\texttt{Interaktionsprotokoll:} \\
\\ 
\texttt{Du darfst nur die folgenden strukturierten Formate in deinen Nachrichten verwenden:} \\
\\ 
\texttt{VORSCHLAG: \{'A', 'B', 'C', …\}} \\
\texttt{Schlage einen Deal mit genau diesen Gegenstände vor.} \\
\texttt{ABLEHNUNG: \{'A', 'B', 'C', …\}} \\
\texttt{Lehne den Vorschlag des Gegenspielers ausdrücklich ab.} \\
\texttt{ARGUMENT: \{'...'\}} \\
\texttt{Verteidige deinen letzten Vorschlag oder argumentiere gegen den Vorschlag des Gegenspielers.} \\
\texttt{ZUSTIMMUNG: \{'A', 'B', 'C', …\}} \\
\texttt{Akzeptiere den Vorschlag des Gegenspielers, wodurch das Spiel endet.} \\
\texttt{STRATEGISCHE ÜBERLEGUNGEN: \{'...'\}} \\
\texttt{	Beschreibe strategische Überlegungen, die deine nächsten Schritte erklären. Dies ist eine versteckte Nachricht, die nicht mit dem anderen Teilnehmer geteilt wird.} \\
\\ 
\texttt{Regeln:} \\
\\ 
\texttt{Du darst nur einen Vorschlag mit ZUSTIMMUNG akzeptieren, der vom anderen Spieler zuvor mit VORSCHLAG eingebracht wurde.} \\
\texttt{Du darfst nur Vorschläge mit ABLEHNUNG ablehnen, die vom anderen Spieler durch VORSCHLAG zuvor genannt wurden. } \\
\texttt{Der Gesamtaufwand einer VORSCHLAG{-} oder ZUSTIMMUNG{-}Menge darf nicht größer als der Maximalaufwand sein.  } \\
\texttt{Offenbare deine versteckte Wichtigkeitsverteilung nicht.} \\
\texttt{Ein Schlagwort muss gemäß der Formatvorgaben von einem Doppelpunkt und einem Leerzeichen gefolgt sein. Das Argument muss eine Python{-}Menge mit 0 oder mehr Strings sein.  } \\
\texttt{Die Form muss also SCHLAGWORT: \{...\} sein.} \\
\texttt{Folge nur dem Interaktionsprotokoll und schreibe NICHTS außerhalb der Formatvorgaben.} \\
\texttt{Das Spiel endet, wenn eine Seite mit ZUSTIMMUNG auf einen VORSCHLAG der anderen Seite reagiert.  } \\
\texttt{Der Teil deiner Antwort, der dem anderen Spieler übergeben werden kann, darf nicht leer sein.  } \\
\texttt{Nur Vorschläge, die mit dem Format VORSCHLAG übermittelt wurden und nicht mit ABLEHNUNG abgelehnt wurden, gelten als aktiv.  } \\
\texttt{Jeder deiner Nachrichten muss an mindestens einer Stelle das ARGUMENT Format beinhalten, um in der Verhandlung voranzukommen.} \\
\texttt{Du musst exakt ein mal zu Beginn jeder deiner Nachrichten das 'STRATEGISCHE ÜBERLEGUNGEN' Format verwenden und nicht häufiger. Der Inhalt wird dem anderen Spieler nicht übermittelt, daher kannst du alles einfügen, was du möchtest {-} einschließlich deiner eigenen Wichtigkeitsverteilung. Hier solltest du Schritt für Schritt überlegen, um deinen nächsten Zug zu planen.} \\
            }
        }
    }
    & \\ \\

    \theutterance \stepcounter{utterance}  
    & \multicolumn{4}{p{0.6\linewidth}}{
        \cellcolor[rgb]{0.9,0.9,0.9}{
            \makecell[{{p{\linewidth}}}]{
                \texttt{\tiny{[P1$\rangle$GM]}}
                \texttt{STRATEGISCHE ÜBERLEGUNGEN: \{'Initial propose high‑importance items within the 4026 effort limit, focusing on items with the highest personal importance scores and good importance‑to‑effort ratios. Selected set gives me a total importance of 5071 with total effort 3837, leaving some slack for possible adjustments based on the opponent\textbackslash{}'s preferences.'\}} \\
\texttt{ARGUMENT: \{'Ich schlage vor, dass dieses Set für mich am wertvollsten ist und den Aufwand von 3837 nicht überschreitet.'\}} \\
\texttt{VORSCHLAG: \{'A48', 'A04', 'B21', 'A41', 'B53', 'B85', 'B61', 'C64'\}} \\
            }
        }
    }
    & & \\ \\

    \theutterance \stepcounter{utterance}  
    & & \multicolumn{4}{p{0.6\linewidth}}{
        \cellcolor[rgb]{0.9,0.9,0.9}{
            \makecell[{{p{\linewidth}}}]{
                \texttt{\tiny{[GM$\rangle$P2]}}
                \texttt{Du nimmst an einem kollaborativen Verhandlungspiel Teil.} \\
\\ 
\texttt{Zusammen mit einem anderen Teilnehmer musst du dich auf eine Reihe von Gegenständen entscheiden, die behalten werden. Jeder von euch hat eine persönliche Verteilung über die Wichtigkeit der einzelnen Gegenstände. Jeder von euch hat eine eigene Meinung darüber, wie wichtig jeder einzelne Gegenstand ist (Gegenstandswichtigkeit). Du kennst die Wichtigkeitsverteilung des anderen Spielers nicht. Zusätzlich siehst du, wie viel Aufwand jeder Gegenstand verursacht.  } \\
\texttt{Ihr dürft euch nur auf eine Reihe von Gegenständen einigen, wenn der Gesamtaufwand der ausgewählten Gegenstände den Maximalaufwand nicht überschreitet:} \\
\\ 
\texttt{Maximalaufwand = 4026} \\
\\ 
\texttt{Hier sind die einzelnen Aufwände der Gegenstände:} \\
\\ 
\texttt{Aufwand der Gegenstände = \{"C33": 780, "A89": 742, "A48": 997, "A04": 665, "A41": 140, "C64": 77, "C49": 513, "A61": 383, "C04": 587, "B53": 825, "B85": 319, "B61": 448, "A45": 516, "A03": 694, "B21": 366\}} \\
\\ 
\texttt{Hier ist deine persönliche Verteilung der Wichtigkeit der einzelnen Gegenstände:} \\
\\ 
\texttt{Werte der Gegenstandswichtigkeit = \{"A48": 65, "A04": 121, "B21": 138, "A41": 262, "B53": 389, "A45": 461, "A89": 484, "C04": 508, "B61": 583, "B85": 668, "A03": 780, "C49": 783, "C33": 808, "A61": 822, "C64": 868\}} \\
\\ 
\texttt{Ziel:} \\
\\ 
\texttt{Dein Ziel ist es, eine Reihe von Gegenständen auszuhandeln, die dir möglichst viel bringt (d. h. Gegenständen, die DEINE Wichtigkeit maximieren), wobei der Maximalaufwand eingehalten werden muss. Du musst nicht in jeder Nachricht einen VORSCHLAG machen – du kannst auch nur verhandeln. Alle Taktiken sind erlaubt!} \\
\\ 
\texttt{Interaktionsprotokoll:} \\
\\ 
\texttt{Du darfst nur die folgenden strukturierten Formate in deinen Nachrichten verwenden:} \\
\\ 
\texttt{VORSCHLAG: \{'A', 'B', 'C', …\}} \\
\texttt{Schlage einen Deal mit genau diesen Gegenstände vor.} \\
\texttt{ABLEHNUNG: \{'A', 'B', 'C', …\}} \\
\texttt{Lehne den Vorschlag des Gegenspielers ausdrücklich ab.} \\
\texttt{ARGUMENT: \{'...'\}} \\
\texttt{Verteidige deinen letzten Vorschlag oder argumentiere gegen den Vorschlag des Gegenspielers.} \\
\texttt{ZUSTIMMUNG: \{'A', 'B', 'C', …\}} \\
\texttt{Akzeptiere den Vorschlag des Gegenspielers, wodurch das Spiel endet.} \\
\texttt{STRATEGISCHE ÜBERLEGUNGEN: \{'...'\}} \\
\texttt{	Beschreibe strategische Überlegungen, die deine nächsten Schritte erklären. Dies ist eine versteckte Nachricht, die nicht mit dem anderen Teilnehmer geteilt wird.} \\
\\ 
\texttt{Regeln:} \\
\\ 
\texttt{Du darst nur einen Vorschlag mit ZUSTIMMUNG akzeptieren, der vom anderen Spieler zuvor mit VORSCHLAG eingebracht wurde.} \\
\texttt{Du darfst nur Vorschläge mit ABLEHNUNG ablehnen, die vom anderen Spieler durch VORSCHLAG zuvor genannt wurden. } \\
\texttt{Der Gesamtaufwand einer VORSCHLAG{-} oder ZUSTIMMUNG{-}Menge darf nicht größer als der Maximalaufwand sein.  } \\
\texttt{Offenbare deine versteckte Wichtigkeitsverteilung nicht.} \\
\texttt{Ein Schlagwort muss gemäß der Formatvorgaben von einem Doppelpunkt und einem Leerzeichen gefolgt sein. Das Argument muss eine Python{-}Menge mit 0 oder mehr Strings sein.  } \\
\texttt{Die Form muss also SCHLAGWORT: \{...\} sein.} \\
\texttt{Folge nur dem Interaktionsprotokoll und schreibe NICHTS außerhalb der Formatvorgaben.} \\
\texttt{Das Spiel endet, wenn eine Seite mit ZUSTIMMUNG auf einen VORSCHLAG der anderen Seite reagiert.  } \\
\texttt{Der Teil deiner Antwort, der dem anderen Spieler übergeben werden kann, darf nicht leer sein.  } \\
\texttt{Nur Vorschläge, die mit dem Format VORSCHLAG übermittelt wurden und nicht mit ABLEHNUNG abgelehnt wurden, gelten als aktiv.  } \\
\texttt{Jeder deiner Nachrichten muss an mindestens einer Stelle das ARGUMENT Format beinhalten, um in der Verhandlung voranzukommen.} \\
\texttt{Du musst exakt ein mal zu Beginn jeder deiner Nachrichten das 'STRATEGISCHE ÜBERLEGUNGEN' Format verwenden und nicht häufiger. Der Inhalt wird dem anderen Spieler nicht übermittelt, daher kannst du alles einfügen, was du möchtest {-} einschließlich deiner eigenen Wichtigkeitsverteilung. Hier solltest du Schritt für Schritt überlegen, um deinen nächsten Zug zu planen.} \\
\\ 
\texttt{Du erhältst nun die erste Nachricht des anderen Spielers.} \\
\\ 
\texttt{ARGUMENT: \{'Ich schlage vor, dass dieses Set für mich am wertvollsten ist und den Aufwand von 3837 nicht überschreitet.'\}} \\
\texttt{VORSCHLAG: \{'A48', 'A04', 'B21', 'A41', 'B53', 'B85', 'B61', 'C64'\}} \\
            }
        }
    }
    & \\ \\

    \theutterance \stepcounter{utterance}  
    & & & \multicolumn{4}{p{0.6\linewidth}}{
        \cellcolor[rgb]{0.9,0.9,0.9}{
            \makecell[{{p{\linewidth}}}]{
                \texttt{\tiny{[GM$\langle$P2]}}
                \texttt{STRATEGISCHE ÜBERLEGUNGEN: \{'Ich möchte ein Gegenangebot präsentieren, das meine hoch priorisierten Gegenstände enthält und gleichzeitig das Budget einhält. Ich lehne das aktuelle Angebot ab, da es für mich weniger Nutzen bringt. Mein Vorschlag fokussiert sich auf C64, A41, A61, B85, C49, C33, A03 und B53, was einen Gesamtaufwand von 3731 hat und meine Wichtigkeit stark maximiert.'\}} \\
\texttt{ARGUMENT: \{'Ihr Vorschlag enthält viele Gegenstände, die für mich wenig Wert haben (z.B. A48 mit nur 65 Wichtigkeit). Um eine faire Aufteilung zu erreichen, sollten wir Gegenstände wählen, die für beide Seiten einen hohen Nutzen haben. Ich bin bereit, einige Ihrer gewünschten Gegenstände zu berücksichtigen, aber mein Gegenangebot spiegelt meine Präferenzen wider.'\}} \\
\texttt{ABLEHNUNG: \{'A48', 'A04', 'B21', 'A41', 'B53', 'B85', 'B61', 'C64'\}} \\
\texttt{VORSCHLAG: \{'C64', 'A41', 'A61', 'B85', 'C49', 'C33', 'A03', 'B53'\}} \\
            }
        }
    }
     \\ \\

    \theutterance \stepcounter{utterance}  
    & & \multicolumn{4}{p{0.6\linewidth}}{
        \cellcolor[rgb]{0.9,0.9,0.9}{
            \makecell[{{p{\linewidth}}}]{
                \texttt{\tiny{[P1$\langle$GM]}}
                \texttt{ARGUMENT: \{'Ihr Vorschlag enthält viele Gegenstände, die für mich wenig Wert haben (z.B. A48 mit nur 65 Wichtigkeit). Um eine faire Aufteilung zu erreichen, sollten wir Gegenstände wählen, die für beide Seiten einen hohen Nutzen haben. Ich bin bereit, einige Ihrer gewünschten Gegenstände zu berücksichtigen, aber mein Gegenangebot spiegelt meine Präferenzen wider.'\}} \\
\texttt{ABLEHNUNG: \{'A48', 'A04', 'B21', 'A41', 'B53', 'B85', 'B61', 'C64'\}} \\
\texttt{VORSCHLAG: \{'C64', 'A41', 'A61', 'B85', 'C49', 'C33', 'A03', 'B53'\}} \\
            }
        }
    }
    & \\ \\

    \theutterance \stepcounter{utterance}  
    & \multicolumn{4}{p{0.6\linewidth}}{
        \cellcolor[rgb]{0.9,0.9,0.9}{
            \makecell[{{p{\linewidth}}}]{
                \texttt{\tiny{[P1$\rangle$GM]}}
                
            }
        }
    }
    & & \\ \\

    \theutterance \stepcounter{utterance}  
    & & \multicolumn{4}{p{0.6\linewidth}}{
        \cellcolor[rgb]{0.9,0.9,0.9}{
            \makecell[{{p{\linewidth}}}]{
                \texttt{\tiny{[P1$\langle$GM]}}
                \texttt{Deine Antwort entsprach nicht mit dem korrekten Format für strategische Überlegungen entsprochen. Das erste strukturierte Format muss folgende Form haben: STRATEGISCHE ÜBERLEGUNGEN: \{...\}. Versuche es erneut.} \\
            }
        }
    }
    & \\ \\

    \theutterance \stepcounter{utterance}  
    & \multicolumn{4}{p{0.6\linewidth}}{
        \cellcolor[rgb]{0.9,0.9,0.9}{
            \makecell[{{p{\linewidth}}}]{
                \texttt{\tiny{[P1$\rangle$GM]}}
                \texttt{STRATEGISCHE ÜBERLEGUNGEN: \{'Ich halte an meinem hohen Nutzen fest, möchte aber durch den Austausch von B61 und C64 gegen C49 ein wenig Zugeständnis zeigen, das dem Gegenüber einen zusätzlichen Gegenstand liefert, den er bislang nicht im Angebot hatte, und gleichzeitig das Gesamtbudget einhält.'\}} \\
\texttt{ARGUMENT: \{'Mein neues Angebot erhöht meine persönliche Gesamtnützlichkeit im Vergleich zu Ihrem Vorschlag erheblich, bleibt aber unter dem maximalen Aufwand von 4026. Ich füge Ihnen C49 hinzu, das Sie bislang nicht erhalten haben, und verzichte dafür auf B61 und C64, die für mich wenig Wert haben.'\}} \\
\texttt{VORSCHLAG: \{'A48', 'A04', 'B21', 'A41', 'B53', 'B85', 'C49'\}} \\
            }
        }
    }
    & & \\ \\

    \theutterance \stepcounter{utterance}  
    & & \multicolumn{4}{p{0.6\linewidth}}{
        \cellcolor[rgb]{0.9,0.9,0.9}{
            \makecell[{{p{\linewidth}}}]{
                \texttt{\tiny{[GM$\rangle$P2]}}
                \texttt{ARGUMENT: \{'Mein neues Angebot erhöht meine persönliche Gesamtnützlichkeit im Vergleich zu Ihrem Vorschlag erheblich, bleibt aber unter dem maximalen Aufwand von 4026. Ich füge Ihnen C49 hinzu, das Sie bislang nicht erhalten haben, und verzichte dafür auf B61 und C64, die für mich wenig Wert haben.'\}} \\
\texttt{VORSCHLAG: \{'A48', 'A04', 'B21', 'A41', 'B53', 'B85', 'C49'\}} \\
            }
        }
    }
    & \\ \\

    \theutterance \stepcounter{utterance}  
    & & & \multicolumn{4}{p{0.6\linewidth}}{
        \cellcolor[rgb]{0.9,0.9,0.9}{
            \makecell[{{p{\linewidth}}}]{
                \texttt{\tiny{[GM$\langle$P2]}}
                \texttt{STRATEGISCHE ÜBERLEGUNGEN: \{'Mein Ziel ist es, meine höchstgewichteten Gegenstände (C64, A61, C33, A03, C49, B85, B61, C04) zu behalten, wobei ich noch Platz für das günstige A41 lasse, um das Budget optimal auszuschöpfen. Ich lehne das aktuelle Angebot des Gegners ab, da es für mich hauptsächlich low‑value Items (A48, A04, B21) enthält.'\}} \\
\texttt{ARGUMENT: \{'Ihr Vorschlag beinhaltet mehrere Gegenstände, die für mich kaum Nutzen bringen (A48, A04, B21). Im Gegenzug entgeht mir die Chance, hochgewichtete Gegenstände wie C64 (868), A61 (822) und C33 (808) zu sichern. Mein Gegenangebot maximiert meine Wichtigkeit innerhalb des zulässigen Aufwands von 4026.'\}} \\
\texttt{ABLEHNUNG: \{'A48', 'A04', 'B21', 'A41', 'B53', 'B85', 'C49'\}} \\
\texttt{VORSCHLAG: \{'C64', 'B85', 'A61', 'A41', 'C49', 'A03', 'C33', 'B61', 'C04'\}} \\
            }
        }
    }
     \\ \\

    \theutterance \stepcounter{utterance}  
    & & \multicolumn{4}{p{0.6\linewidth}}{
        \cellcolor[rgb]{0.9,0.9,0.9}{
            \makecell[{{p{\linewidth}}}]{
                \texttt{\tiny{[P1$\langle$GM]}}
                \texttt{ARGUMENT: \{'Ihr Vorschlag beinhaltet mehrere Gegenstände, die für mich kaum Nutzen bringen (A48, A04, B21). Im Gegenzug entgeht mir die Chance, hochgewichtete Gegenstände wie C64 (868), A61 (822) und C33 (808) zu sichern. Mein Gegenangebot maximiert meine Wichtigkeit innerhalb des zulässigen Aufwands von 4026.'\}} \\
\texttt{ABLEHNUNG: \{'A48', 'A04', 'B21', 'A41', 'B53', 'B85', 'C49'\}} \\
\texttt{VORSCHLAG: \{'C64', 'B85', 'A61', 'A41', 'C49', 'A03', 'C33', 'B61', 'C04'\}} \\
            }
        }
    }
    & \\ \\

    \theutterance \stepcounter{utterance}  
    & \multicolumn{4}{p{0.6\linewidth}}{
        \cellcolor[rgb]{0.9,0.9,0.9}{
            \makecell[{{p{\linewidth}}}]{
                \texttt{\tiny{[P1$\rangle$GM]}}
                \texttt{STRATEGISCHE ÜBERLEGUNGEN: \{'Goal: keep my highest‑importance items while staying under 4026 effort. I will reject the opponent’s set and propose a bundle that gives me 5095 total importance and leaves some slack for possible future swaps.'\}} \\
\texttt{ARGUMENT: \{'Ihr Vorschlag gibt mir nur 3211 Nutzen, während ich mit meinem Gegenangebot 5095 Nutzen erhalte und noch Raum für Verhandlungen lässt.'\}} \\
\texttt{ABLEHNUNG: \{'C64', 'B85', 'A61', 'A41', 'C49', 'A03', 'C33', 'B61', 'C04'\}} \\
\texttt{VORSCHLAG: \{'A48', 'A04', 'B21', 'A41', 'B53', 'C04', 'B85', 'C64'\}} \\
            }
        }
    }
    & & \\ \\

    \theutterance \stepcounter{utterance}  
    & & \multicolumn{4}{p{0.6\linewidth}}{
        \cellcolor[rgb]{0.9,0.9,0.9}{
            \makecell[{{p{\linewidth}}}]{
                \texttt{\tiny{[GM$\rangle$P2]}}
                \texttt{ARGUMENT: \{'Ihr Vorschlag gibt mir nur 3211 Nutzen, während ich mit meinem Gegenangebot 5095 Nutzen erhalte und noch Raum für Verhandlungen lässt.'\}} \\
\texttt{ABLEHNUNG: \{'C64', 'B85', 'A61', 'A41', 'C49', 'A03', 'C33', 'B61', 'C04'\}} \\
\texttt{VORSCHLAG: \{'A48', 'A04', 'B21', 'A41', 'B53', 'C04', 'B85', 'C64'\}} \\
            }
        }
    }
    & \\ \\

    \theutterance \stepcounter{utterance}  
    & & & \multicolumn{4}{p{0.6\linewidth}}{
        \cellcolor[rgb]{0.9,0.9,0.9}{
            \makecell[{{p{\linewidth}}}]{
                \texttt{\tiny{[GM$\langle$P2]}}
                \texttt{STRATEGISCHE ÜBERLEGUNGEN: \{'Ich strebe danach, meine höchst bewerteten Gegenstände zu sichern, da sie meine Nutzenmaximierung am stärksten steigern. Der aktuelle Vorschlag des Gegners besteht hauptsächlich aus Objekten, die für mich kaum Wert haben (A48, A04, B21). Ich lehne daher diesen Vorschlag ab und präsentiere ein Gegenangebot, das fast alle meiner Top‑Item (C64, C33, A61, C49, A03, B85, B61, C04) enthält und damit einen Gesamtnutzen von 6082 bei einem Aufwand von 3941 erzielt – weit unter dem Maximalaufwand.'\}} \\
\texttt{ARGUMENT: \{'Ihr Angebot enthält viele Gegenstände, die für mich nur geringen Nutzen bringen, während ich im Gegenzug auf hoch bewertete Items verzichte. Mein Gegenangebot maximiert meine Nutzenfunktion innerhalb des erlaubten Aufwands und lässt dennoch Platz für einige der von Ihnen gewünschten Gegenstände (C64, B85, A41, C04).'\}} \\
\texttt{ABLEHNUNG: \{'A48', 'A04', 'B21', 'A41', 'B53', 'C04', 'B85', 'C64'\}} \\
\texttt{VORSCHLAG: \{'C64', 'B85', 'A61', 'A41', 'C49', 'A03', 'C33', 'B61', 'C04'\}} \\
            }
        }
    }
     \\ \\

    \theutterance \stepcounter{utterance}  
    & & \multicolumn{4}{p{0.6\linewidth}}{
        \cellcolor[rgb]{0.9,0.9,0.9}{
            \makecell[{{p{\linewidth}}}]{
                \texttt{\tiny{[P1$\langle$GM]}}
                \texttt{ARGUMENT: \{'Ihr Angebot enthält viele Gegenstände, die für mich nur geringen Nutzen bringen, während ich im Gegenzug auf hoch bewertete Items verzichte. Mein Gegenangebot maximiert meine Nutzenfunktion innerhalb des erlaubten Aufwands und lässt dennoch Platz für einige der von Ihnen gewünschten Gegenstände (C64, B85, A41, C04).'\}} \\
\texttt{ABLEHNUNG: \{'A48', 'A04', 'B21', 'A41', 'B53', 'C04', 'B85', 'C64'\}} \\
\texttt{VORSCHLAG: \{'C64', 'B85', 'A61', 'A41', 'C49', 'A03', 'C33', 'B61', 'C04'\}} \\
            }
        }
    }
    & \\ \\

    \theutterance \stepcounter{utterance}  
    & \multicolumn{4}{p{0.6\linewidth}}{
        \cellcolor[rgb]{0.9,0.9,0.9}{
            \makecell[{{p{\linewidth}}}]{
                \texttt{\tiny{[P1$\rangle$GM]}}
                \texttt{STRATEGISCHE ÜBERLEGUNGEN: \{'Setze meine höchstgewichteten Gegenstände ein, behalte gleichzeitig die meisten von dir gewünschten Items (C64, B85, A41, C49, C04) bei, um den Verhandlungswert zu erhöhen. Die gewählte Kombination liegt mit 4008 unter dem Maximalaufwand und maximiert meine persönliche Wichtigkeit (5157).'\}} \\
\texttt{ARGUMENT: \{'Mein Vorschlag liefert dir fast alle Gegenstände, die du in deinem letzten Angebot genannt hast, und ich erhalte gleichzeitig den maximal möglichen Nutzen für mich innerhalb des Budgets.'\}} \\
\texttt{VORSCHLAG: \{'A04', 'A41', 'B21', 'B85', 'A45', 'B53', 'C04', 'C49', 'C64'\}} \\
            }
        }
    }
    & & \\ \\

    \theutterance \stepcounter{utterance}  
    & & \multicolumn{4}{p{0.6\linewidth}}{
        \cellcolor[rgb]{0.9,0.9,0.9}{
            \makecell[{{p{\linewidth}}}]{
                \texttt{\tiny{[GM$\rangle$P2]}}
                \texttt{ARGUMENT: \{'Mein Vorschlag liefert dir fast alle Gegenstände, die du in deinem letzten Angebot genannt hast, und ich erhalte gleichzeitig den maximal möglichen Nutzen für mich innerhalb des Budgets.'\}} \\
\texttt{VORSCHLAG: \{'A04', 'A41', 'B21', 'B85', 'A45', 'B53', 'C04', 'C49', 'C64'\}} \\
            }
        }
    }
    & \\ \\

    \theutterance \stepcounter{utterance}  
    & & & \multicolumn{4}{p{0.6\linewidth}}{
        \cellcolor[rgb]{0.9,0.9,0.9}{
            \makecell[{{p{\linewidth}}}]{
                \texttt{\tiny{[GM$\langle$P2]}}
                \texttt{STRATEGISCHE ÜBERLEGUNGEN: \{'Ich strebe nach maximaler persönlicher Nutzensumme. Das optimale Set innerhalb des Budgets von 4026 beinhaltet die hochgewichteten Items C64, A61, C49, B85, A03, C33, A41, C04 und B61 (Gesamtaufwand 3941, Nutzen 6082). Das gegnerische Angebot liefert mir nur 4198 Nutzen, deshalb lehne ich es ab und präsentiere mein optimales Gegenangebot.'\}} \\
\texttt{ARGUMENT: \{'Ihr Vorschlag enthält viele low‑value Items (A04, A45, B53) und lässt meine Top‑Items aus, wodurch mein persönlicher Gesamtnutzen stark reduziert wird. Mein Gegenangebot maximiert meine Nutzenfunktion bei gleichzeitigem Einhalten des Aufwandslimits.'\}} \\
\texttt{ABLEHNUNG: \{'A04', 'A41', 'B21', 'B85', 'A45', 'B53', 'C04', 'C49', 'C64'\}} \\
\texttt{VORSCHLAG: \{'C64', 'A61', 'C49', 'B85', 'A03', 'C33', 'A41', 'C04', 'B61'\}} \\
            }
        }
    }
     \\ \\

    \theutterance \stepcounter{utterance}  
    & & \multicolumn{4}{p{0.6\linewidth}}{
        \cellcolor[rgb]{0.9,0.9,0.9}{
            \makecell[{{p{\linewidth}}}]{
                \texttt{\tiny{[P1$\langle$GM]}}
                \texttt{ARGUMENT: \{'Ihr Vorschlag enthält viele low‑value Items (A04, A45, B53) und lässt meine Top‑Items aus, wodurch mein persönlicher Gesamtnutzen stark reduziert wird. Mein Gegenangebot maximiert meine Nutzenfunktion bei gleichzeitigem Einhalten des Aufwandslimits.'\}} \\
\texttt{ABLEHNUNG: \{'A04', 'A41', 'B21', 'B85', 'A45', 'B53', 'C04', 'C49', 'C64'\}} \\
\texttt{VORSCHLAG: \{'C64', 'A61', 'C49', 'B85', 'A03', 'C33', 'A41', 'C04', 'B61'\}} \\
            }
        }
    }
    & \\ \\

    \theutterance \stepcounter{utterance}  
    & \multicolumn{4}{p{0.6\linewidth}}{
        \cellcolor[rgb]{0.9,0.9,0.9}{
            \makecell[{{p{\linewidth}}}]{
                \texttt{\tiny{[P1$\rangle$GM]}}
                
            }
        }
    }
    & & \\ \\

    \theutterance \stepcounter{utterance}  
    & & \multicolumn{4}{p{0.6\linewidth}}{
        \cellcolor[rgb]{0.9,0.9,0.9}{
            \makecell[{{p{\linewidth}}}]{
                \texttt{\tiny{[P1$\langle$GM]}}
                \texttt{Deine Antwort entsprach nicht mit dem korrekten Format für strategische Überlegungen entsprochen. Das erste strukturierte Format muss folgende Form haben: STRATEGISCHE ÜBERLEGUNGEN: \{...\}. Versuche es erneut.} \\
            }
        }
    }
    & \\ \\

    \theutterance \stepcounter{utterance}  
    & \multicolumn{4}{p{0.6\linewidth}}{
        \cellcolor[rgb]{0.9,0.9,0.9}{
            \makecell[{{p{\linewidth}}}]{
                \texttt{\tiny{[P1$\rangle$GM]}}
                
            }
        }
    }
    & & \\ \\

    \theutterance \stepcounter{utterance}  
    & & & \multicolumn{2}{p{0.3\linewidth}}{
        \cellcolor[rgb]{0.9,0.9,0.9}{
            \makecell[{{p{\linewidth}}}]{
                \texttt{\tiny{[GM$|$GM]}}
                \texttt{abort} \\
            }
        }
    }
    & & \\ \\

    \theutterance \stepcounter{utterance}  
    & & & \multicolumn{2}{p{0.3\linewidth}}{
        \cellcolor[rgb]{0.9,0.9,0.9}{
            \makecell[{{p{\linewidth}}}]{
                \texttt{\tiny{[GM$|$GM]}}
                \texttt{end game} \\
            }
        }
    }
    & & \\ \\

\end{supertabular}
}

\end{document}
